% 分野別類題: クレロー方程式・ラグランジュ方程式（1階・非正規形）
% p = y' とおいて微分し、一般解（直線族）と特異解（包絡線）を求める。
% コンパイル: TEXINPUTS="../../../00_common:$TEXINPUTS" latexmk -xelatex problems.tex
\documentclass[a4paper,11pt]{article}
\usepackage{preamble}

\begin{document}

\kamokumei{類題演習 --- クレロー方程式・ラグランジュ方程式}

\shijibun{次の常微分方程式を解き、導出過程を示せ。（クレロー方程式 $y = xy' + f(y')$ は $p = y'$ とおいて両辺を $x$ で微分すると $(x + f'(p))\dfrac{dp}{dx} = 0$ となり、一般解（直線族）$y = Cx + f(C)$ と、特異解（$p$ を消去して得られる一般解の包絡線）が得られる。ラグランジュ方程式 $y = xg(y') + f(y')$（$g(p) \not\equiv p$）は同様に $p = y'$ とおき、$x$ を $p$ の関数とみなして両辺を $p$ で微分することで、$x$ に関する1階線形微分方程式に帰着させて解く。）}

\shomon{問1.}{次のクレロー方程式の一般解と特異解を求めよ。}
\[
y = xy' + (y')^{2}
\]

\shomon{問2.}{次のクレロー方程式の一般解と特異解を求めよ（$y' > 0$ とする）。}
\[
y = xy' - \ln y'
\]

\shomon{問3.}{次のラグランジュ方程式について、$p = y'$ をパラメータとして一般解 $x = x(p)$, $y = y(p)$ を求めよ。さらに、傾き $p = 1$ のとき曲線が点 $\left(x, y\right) = \left(-\dfrac{1}{3}, \dfrac{1}{3}\right)$ を通るように、任意定数 $C$ の値を定めよ。}
\[
y = 2xy' + (y')^{2}
\]

\end{document}
